\documentclass[11pt]{article}
\newcommand{\videoid}{Sot3oC0gpmI}
% Preámbulo compartido de las notas del curso de topología algebraica.
% Cada nota empieza con \documentclass[11pt]{article}, define \videoid{...}
% (el identificador del video de YouTube) y luego hace % Preámbulo compartido de las notas del curso de topología algebraica.
% Cada nota empieza con \documentclass[11pt]{article}, define \videoid{...}
% (el identificador del video de YouTube) y luego hace % Preámbulo compartido de las notas del curso de topología algebraica.
% Cada nota empieza con \documentclass[11pt]{article}, define \videoid{...}
% (el identificador del video de YouTube) y luego hace \input{preambulo}.
% Compilar con LuaLaTeX:  lualatex nombre-de-la-nota.tex
\usepackage[spanish,provide=*]{babel}
\usepackage[T1]{fontenc}
\usepackage{lmodern}
\usepackage[a4paper,margin=2.6cm]{geometry}
\usepackage{amsmath,amssymb,amsthm}
\usepackage{xcolor}
\usepackage{tikz}
\usetikzlibrary{arrows.meta,babel}
\usepackage[colorlinks=true,urlcolor=blue!50!black,linkcolor=blue!50!black]{hyperref}

\hyphenpenalty=10000 \exhyphenpenalty=10000 \emergencystretch=3em

% Colores de las figuras (los mismos que usa la página)
\definecolor{extremo}{RGB}{47,107,79}
\definecolor{azul}{RGB}{44,74,154}

% Definiciones numeradas por unidad
\theoremstyle{definition}
\newtheorem{defn}{Definición}
% Numeración por unidad: \unidad{u}{n} (después de \input{preambulo}) hace que las
% definiciones de la nota se llamen u.n, u.(n+1), ...; n es el primer número
% libre de esa unidad (ver la tabla de CONVENCIONES.md).
\newcommand{\unidad}[2]{\renewcommand{\thedefn}{#1.\arabic{defn}}\setcounter{defn}{\numexpr#2-1\relax}}

% Observaciones y erratas del video. Usar SOLO si hay algo que señalar:
%   \begin{erratas}
%     \errata{203}{3:23}{Se dice ..., pero debe decir ...}
%   \end{erratas}
% El primer argumento son los segundos (enlaza al minuto exacto del video),
% el segundo es la marca visible, el tercero es el texto.
\newenvironment{erratas}{\section*{Observaciones y erratas del video}\begin{itemize}}{\end{itemize}}
\newcommand{\errata}[3]{%
  \edef\errurl{https://www.youtube.com/watch?v=\videoid\string&t=#1}%
  \item[\textbf{\expandafter\href\expandafter{\errurl}{#2}}] #3}
.
% Compilar con LuaLaTeX:  lualatex nombre-de-la-nota.tex
\usepackage[spanish,provide=*]{babel}
\usepackage[T1]{fontenc}
\usepackage{lmodern}
\usepackage[a4paper,margin=2.6cm]{geometry}
\usepackage{amsmath,amssymb,amsthm}
\usepackage{xcolor}
\usepackage{tikz}
\usetikzlibrary{arrows.meta,babel}
\usepackage[colorlinks=true,urlcolor=blue!50!black,linkcolor=blue!50!black]{hyperref}

\hyphenpenalty=10000 \exhyphenpenalty=10000 \emergencystretch=3em

% Colores de las figuras (los mismos que usa la página)
\definecolor{extremo}{RGB}{47,107,79}
\definecolor{azul}{RGB}{44,74,154}

% Definiciones numeradas por unidad
\theoremstyle{definition}
\newtheorem{defn}{Definición}
% Numeración por unidad: \unidad{u}{n} (después de % Preámbulo compartido de las notas del curso de topología algebraica.
% Cada nota empieza con \documentclass[11pt]{article}, define \videoid{...}
% (el identificador del video de YouTube) y luego hace \input{preambulo}.
% Compilar con LuaLaTeX:  lualatex nombre-de-la-nota.tex
\usepackage[spanish,provide=*]{babel}
\usepackage[T1]{fontenc}
\usepackage{lmodern}
\usepackage[a4paper,margin=2.6cm]{geometry}
\usepackage{amsmath,amssymb,amsthm}
\usepackage{xcolor}
\usepackage{tikz}
\usetikzlibrary{arrows.meta,babel}
\usepackage[colorlinks=true,urlcolor=blue!50!black,linkcolor=blue!50!black]{hyperref}

\hyphenpenalty=10000 \exhyphenpenalty=10000 \emergencystretch=3em

% Colores de las figuras (los mismos que usa la página)
\definecolor{extremo}{RGB}{47,107,79}
\definecolor{azul}{RGB}{44,74,154}

% Definiciones numeradas por unidad
\theoremstyle{definition}
\newtheorem{defn}{Definición}
% Numeración por unidad: \unidad{u}{n} (después de \input{preambulo}) hace que las
% definiciones de la nota se llamen u.n, u.(n+1), ...; n es el primer número
% libre de esa unidad (ver la tabla de CONVENCIONES.md).
\newcommand{\unidad}[2]{\renewcommand{\thedefn}{#1.\arabic{defn}}\setcounter{defn}{\numexpr#2-1\relax}}

% Observaciones y erratas del video. Usar SOLO si hay algo que señalar:
%   \begin{erratas}
%     \errata{203}{3:23}{Se dice ..., pero debe decir ...}
%   \end{erratas}
% El primer argumento son los segundos (enlaza al minuto exacto del video),
% el segundo es la marca visible, el tercero es el texto.
\newenvironment{erratas}{\section*{Observaciones y erratas del video}\begin{itemize}}{\end{itemize}}
\newcommand{\errata}[3]{%
  \edef\errurl{https://www.youtube.com/watch?v=\videoid\string&t=#1}%
  \item[\textbf{\expandafter\href\expandafter{\errurl}{#2}}] #3}
) hace que las
% definiciones de la nota se llamen u.n, u.(n+1), ...; n es el primer número
% libre de esa unidad (ver la tabla de CONVENCIONES.md).
\newcommand{\unidad}[2]{\renewcommand{\thedefn}{#1.\arabic{defn}}\setcounter{defn}{\numexpr#2-1\relax}}

% Observaciones y erratas del video. Usar SOLO si hay algo que señalar:
%   \begin{erratas}
%     \errata{203}{3:23}{Se dice ..., pero debe decir ...}
%   \end{erratas}
% El primer argumento son los segundos (enlaza al minuto exacto del video),
% el segundo es la marca visible, el tercero es el texto.
\newenvironment{erratas}{\section*{Observaciones y erratas del video}\begin{itemize}}{\end{itemize}}
\newcommand{\errata}[3]{%
  \edef\errurl{https://www.youtube.com/watch?v=\videoid\string&t=#1}%
  \item[\textbf{\expandafter\href\expandafter{\errurl}{#2}}] #3}
.
% Compilar con LuaLaTeX:  lualatex nombre-de-la-nota.tex
\usepackage[spanish,provide=*]{babel}
\usepackage[T1]{fontenc}
\usepackage{lmodern}
\usepackage[a4paper,margin=2.6cm]{geometry}
\usepackage{amsmath,amssymb,amsthm}
\usepackage{xcolor}
\usepackage{tikz}
\usetikzlibrary{arrows.meta,babel}
\usepackage[colorlinks=true,urlcolor=blue!50!black,linkcolor=blue!50!black]{hyperref}

\hyphenpenalty=10000 \exhyphenpenalty=10000 \emergencystretch=3em

% Colores de las figuras (los mismos que usa la página)
\definecolor{extremo}{RGB}{47,107,79}
\definecolor{azul}{RGB}{44,74,154}

% Definiciones numeradas por unidad
\theoremstyle{definition}
\newtheorem{defn}{Definición}
% Numeración por unidad: \unidad{u}{n} (después de % Preámbulo compartido de las notas del curso de topología algebraica.
% Cada nota empieza con \documentclass[11pt]{article}, define \videoid{...}
% (el identificador del video de YouTube) y luego hace % Preámbulo compartido de las notas del curso de topología algebraica.
% Cada nota empieza con \documentclass[11pt]{article}, define \videoid{...}
% (el identificador del video de YouTube) y luego hace \input{preambulo}.
% Compilar con LuaLaTeX:  lualatex nombre-de-la-nota.tex
\usepackage[spanish,provide=*]{babel}
\usepackage[T1]{fontenc}
\usepackage{lmodern}
\usepackage[a4paper,margin=2.6cm]{geometry}
\usepackage{amsmath,amssymb,amsthm}
\usepackage{xcolor}
\usepackage{tikz}
\usetikzlibrary{arrows.meta,babel}
\usepackage[colorlinks=true,urlcolor=blue!50!black,linkcolor=blue!50!black]{hyperref}

\hyphenpenalty=10000 \exhyphenpenalty=10000 \emergencystretch=3em

% Colores de las figuras (los mismos que usa la página)
\definecolor{extremo}{RGB}{47,107,79}
\definecolor{azul}{RGB}{44,74,154}

% Definiciones numeradas por unidad
\theoremstyle{definition}
\newtheorem{defn}{Definición}
% Numeración por unidad: \unidad{u}{n} (después de \input{preambulo}) hace que las
% definiciones de la nota se llamen u.n, u.(n+1), ...; n es el primer número
% libre de esa unidad (ver la tabla de CONVENCIONES.md).
\newcommand{\unidad}[2]{\renewcommand{\thedefn}{#1.\arabic{defn}}\setcounter{defn}{\numexpr#2-1\relax}}

% Observaciones y erratas del video. Usar SOLO si hay algo que señalar:
%   \begin{erratas}
%     \errata{203}{3:23}{Se dice ..., pero debe decir ...}
%   \end{erratas}
% El primer argumento son los segundos (enlaza al minuto exacto del video),
% el segundo es la marca visible, el tercero es el texto.
\newenvironment{erratas}{\section*{Observaciones y erratas del video}\begin{itemize}}{\end{itemize}}
\newcommand{\errata}[3]{%
  \edef\errurl{https://www.youtube.com/watch?v=\videoid\string&t=#1}%
  \item[\textbf{\expandafter\href\expandafter{\errurl}{#2}}] #3}
.
% Compilar con LuaLaTeX:  lualatex nombre-de-la-nota.tex
\usepackage[spanish,provide=*]{babel}
\usepackage[T1]{fontenc}
\usepackage{lmodern}
\usepackage[a4paper,margin=2.6cm]{geometry}
\usepackage{amsmath,amssymb,amsthm}
\usepackage{xcolor}
\usepackage{tikz}
\usetikzlibrary{arrows.meta,babel}
\usepackage[colorlinks=true,urlcolor=blue!50!black,linkcolor=blue!50!black]{hyperref}

\hyphenpenalty=10000 \exhyphenpenalty=10000 \emergencystretch=3em

% Colores de las figuras (los mismos que usa la página)
\definecolor{extremo}{RGB}{47,107,79}
\definecolor{azul}{RGB}{44,74,154}

% Definiciones numeradas por unidad
\theoremstyle{definition}
\newtheorem{defn}{Definición}
% Numeración por unidad: \unidad{u}{n} (después de % Preámbulo compartido de las notas del curso de topología algebraica.
% Cada nota empieza con \documentclass[11pt]{article}, define \videoid{...}
% (el identificador del video de YouTube) y luego hace \input{preambulo}.
% Compilar con LuaLaTeX:  lualatex nombre-de-la-nota.tex
\usepackage[spanish,provide=*]{babel}
\usepackage[T1]{fontenc}
\usepackage{lmodern}
\usepackage[a4paper,margin=2.6cm]{geometry}
\usepackage{amsmath,amssymb,amsthm}
\usepackage{xcolor}
\usepackage{tikz}
\usetikzlibrary{arrows.meta,babel}
\usepackage[colorlinks=true,urlcolor=blue!50!black,linkcolor=blue!50!black]{hyperref}

\hyphenpenalty=10000 \exhyphenpenalty=10000 \emergencystretch=3em

% Colores de las figuras (los mismos que usa la página)
\definecolor{extremo}{RGB}{47,107,79}
\definecolor{azul}{RGB}{44,74,154}

% Definiciones numeradas por unidad
\theoremstyle{definition}
\newtheorem{defn}{Definición}
% Numeración por unidad: \unidad{u}{n} (después de \input{preambulo}) hace que las
% definiciones de la nota se llamen u.n, u.(n+1), ...; n es el primer número
% libre de esa unidad (ver la tabla de CONVENCIONES.md).
\newcommand{\unidad}[2]{\renewcommand{\thedefn}{#1.\arabic{defn}}\setcounter{defn}{\numexpr#2-1\relax}}

% Observaciones y erratas del video. Usar SOLO si hay algo que señalar:
%   \begin{erratas}
%     \errata{203}{3:23}{Se dice ..., pero debe decir ...}
%   \end{erratas}
% El primer argumento son los segundos (enlaza al minuto exacto del video),
% el segundo es la marca visible, el tercero es el texto.
\newenvironment{erratas}{\section*{Observaciones y erratas del video}\begin{itemize}}{\end{itemize}}
\newcommand{\errata}[3]{%
  \edef\errurl{https://www.youtube.com/watch?v=\videoid\string&t=#1}%
  \item[\textbf{\expandafter\href\expandafter{\errurl}{#2}}] #3}
) hace que las
% definiciones de la nota se llamen u.n, u.(n+1), ...; n es el primer número
% libre de esa unidad (ver la tabla de CONVENCIONES.md).
\newcommand{\unidad}[2]{\renewcommand{\thedefn}{#1.\arabic{defn}}\setcounter{defn}{\numexpr#2-1\relax}}

% Observaciones y erratas del video. Usar SOLO si hay algo que señalar:
%   \begin{erratas}
%     \errata{203}{3:23}{Se dice ..., pero debe decir ...}
%   \end{erratas}
% El primer argumento son los segundos (enlaza al minuto exacto del video),
% el segundo es la marca visible, el tercero es el texto.
\newenvironment{erratas}{\section*{Observaciones y erratas del video}\begin{itemize}}{\end{itemize}}
\newcommand{\errata}[3]{%
  \edef\errurl{https://www.youtube.com/watch?v=\videoid\string&t=#1}%
  \item[\textbf{\expandafter\href\expandafter{\errurl}{#2}}] #3}
) hace que las
% definiciones de la nota se llamen u.n, u.(n+1), ...; n es el primer número
% libre de esa unidad (ver la tabla de CONVENCIONES.md).
\newcommand{\unidad}[2]{\renewcommand{\thedefn}{#1.\arabic{defn}}\setcounter{defn}{\numexpr#2-1\relax}}

% Observaciones y erratas del video. Usar SOLO si hay algo que señalar:
%   \begin{erratas}
%     \errata{203}{3:23}{Se dice ..., pero debe decir ...}
%   \end{erratas}
% El primer argumento son los segundos (enlaza al minuto exacto del video),
% el segundo es la marca visible, el tercero es el texto.
\newenvironment{erratas}{\section*{Observaciones y erratas del video}\begin{itemize}}{\end{itemize}}
\newcommand{\errata}[3]{%
  \edef\errurl{https://www.youtube.com/watch?v=\videoid\string&t=#1}%
  \item[\textbf{\expandafter\href\expandafter{\errurl}{#2}}] #3}

\unidad{4}{10} % unidad 4; primer enunciado: 4.10
\definecolor{serpiente}{RGB}{176,64,150}

\title{Álgebra homológica: el lema de la serpiente}
\author{Luis Jorge Sánchez Saldaña\\[2pt]\small Facultad de Ciencias, UNAM}
\date{}

\begin{document}
\maketitle

\noindent\small Notas del video
\href{https://www.youtube.com/watch?v=Sot3oC0gpmI}{\guillemotleft{}Álgebra homológica - el lema de la serpiente\guillemotright{}}
(curso de Introducción a la Topología Algebraica). Cuarto video de la serie de álgebra homológica: se enuncia y se demuestra el lema de la serpiente, que se usará para demostrar el teorema fundamental del álgebra homológica.
\normalsize

\section*{El enunciado}

El lema de la serpiente servirá como auxiliar para demostrar, posteriormente, el teorema fundamental del álgebra homológica.

\begin{lema}[Lema de la serpiente]\label{lema:serpiente}
Considérese el diagrama conmutativo de $R$-módulos
\[
  \begin{array}{ccccccccc}
    & & Z'_1 & \xrightarrow{\ \phi'\ } & Z'_2 & \xrightarrow{\ \psi'\ } & Z'_3 & \longrightarrow & 0\\[2pt]
    & & {\scriptstyle\partial_1}\big\downarrow & & \big\downarrow{\scriptstyle\partial_2} & & \big\downarrow{\scriptstyle\partial_3} & & \\[2pt]
    0 & \longrightarrow & Z_1 & \xrightarrow{\ \phi\ } & Z_2 & \xrightarrow{\ \psi\ } & Z_3 & &
  \end{array}
\]
y supóngase que las filas son exactas. Entonces existe un homomorfismo
\[
  \delta_*\colon\ker(\partial_3)\longrightarrow Z_1/\operatorname{Im}(\partial_1)
\]
tal que la sucesión
\[
  \ker(\partial_1)\xrightarrow{\ \phi''\ }\ker(\partial_2)\xrightarrow{\ \psi''\ }\ker(\partial_3)\xrightarrow{\ \delta_*\ }Z_1/\operatorname{Im}(\partial_1)\xrightarrow{\ \bar\phi\ }Z_2/\operatorname{Im}(\partial_2)\xrightarrow{\ \bar\psi\ }Z_3/\operatorname{Im}(\partial_3)
\]
es exacta, donde $\phi''$ y $\psi''$ son las restricciones de $\phi'$ y $\psi'$, y $\bar\phi$, $\bar\psi$ son los homomorfismos inducidos por $\phi$ y $\psi$.
\end{lema}

Es decir, las funciones $\phi''$, $\psi''$, $\bar\phi$ y $\bar\psi$ son las naturales: las dos primeras son restricciones de $\phi'$ y $\psi'$, y las otras dos son las inducidas por $\phi$ y $\psi$ en los cocientes.

\medskip
La pregunta natural es: ?`dónde está la serpiente? Hay que poner el diagrama de manera adecuada. Arriba del diagrama se ponen los núcleos de $\partial_1$, $\partial_2$ y $\partial_3$ con sus inclusiones, y abajo los cocientes de $Z_1$, $Z_2$ y $Z_3$ entre las imágenes de $\partial_1$, $\partial_2$ y $\partial_3$, con las proyecciones al cociente. Entre los núcleos están las restricciones $\phi''$ y $\psi''$, y entre los cocientes, las inducidas $\bar\phi$ y $\bar\psi$. Ahí aparece la serpiente: es el homomorfismo $\delta_*$, que sale de $\ker\partial_3$, atraviesa todo el diagrama y llega a $Z_1/\operatorname{Im}\partial_1$. El enunciado dice que la sucesión que comienza en $\ker\partial_1$, se conecta mediante la serpiente y continúa hasta $Z_3/\operatorname{Im}\partial_3$ es exacta.

\begin{center}
\begin{tikzpicture}[x=1cm,y=1.3cm,>=Stealth]
  \node (k1) at (2.8,3) {$\ker\partial_1$}; \node (k2) at (5.6,3) {$\ker\partial_2$}; \node (k3) at (8.4,3) {$\ker\partial_3$};
  \node (a1) at (2.8,2) {$Z'_1$}; \node (a2) at (5.6,2) {$Z'_2$}; \node (a3) at (8.4,2) {$Z'_3$}; \node (a4) at (10.6,2) {$0$};
  \node (b0) at (0.6,1) {$0$}; \node (b1) at (2.8,1) {$Z_1$}; \node (b2) at (5.6,1) {$Z_2$}; \node (b3) at (8.4,1) {$Z_3$};
  \node (c1) at (2.8,0) {$Z_1/\operatorname{Im}\partial_1$}; \node (c2) at (5.6,0) {$Z_2/\operatorname{Im}\partial_2$}; \node (c3) at (8.4,0) {$Z_3/\operatorname{Im}\partial_3$};
  \draw[->] (k1) -- node[above,font=\small] {$\phi''$} (k2); \draw[->] (k2) -- node[above,font=\small] {$\psi''$} (k3);
  \draw[->] (a1) -- node[above,font=\small] {$\phi'$} (a2); \draw[->] (a2) -- node[above,font=\small] {$\psi'$} (a3); \draw[->] (a3) -- (a4);
  \draw[->] (b0) -- (b1); \draw[->] (b1) -- node[below,font=\small] {$\phi$} (b2); \draw[->] (b2) -- node[below,font=\small] {$\psi$} (b3);
  \draw[->] (c1) -- node[above,font=\small] {$\bar\phi$} (c2); \draw[->] (c2) -- node[above,font=\small] {$\bar\psi$} (c3);
  \foreach \i in {1,2,3} {
    \draw[{Hooks[right]}->] (k\i) -- (a\i);
    \draw[->] (a\i) -- node[left,font=\small,pos=.25] {$\partial_\i$} (b\i);
    \draw[->>] (b\i) -- (c\i);
  }
  \draw[->,serpiente,thick,rounded corners=8pt] (k3.east) -- (9.9,3) -- (9.9,1.5) -- (-0.5,1.5) -- (-0.5,0) -- (c1.west);
  \node[serpiente,font=\small,fill=white,inner sep=1pt] at (4.2,1.5) {$\delta_*$};
\end{tikzpicture}
\end{center}

\section*{La demostración}

Hay que demostrar la exactitud en cuatro lugares: en $\ker\partial_2$, en $\ker\partial_3$, en $Z_1/\operatorname{Im}\partial_1$ y en $Z_2/\operatorname{Im}\partial_2$. Es decir, en cada uno de esos módulos, que la imagen del homomorfismo de la izquierda es igual al núcleo del homomorfismo de la derecha. Pero lo primero es demostrar que $\phi''$, $\psi''$, $\bar\phi$ y $\bar\psi$ están bien definidas.

\paragraph{$\phi''$ está bien definida.} Sea $c\in\ker\partial_1$; se quiere ver que $\phi''(c)\in\ker\partial_2$. Aplicando $\partial_2$ y usando la conmutatividad del diagrama,
\[
  \partial_2(\phi'(c))=\phi(\partial_1(c))=\phi(0)=0 .
\]
Análogamente, $\psi''$ está bien definida.

\paragraph{$\bar\phi$ está bien definida.} Basta ver que $\phi(\operatorname{Im}\partial_1)\subseteq\operatorname{Im}\partial_2$, y entonces se tiene la función inducida en los cocientes correspondientes. Sea $x\in\operatorname{Im}\partial_1$, es decir, $x=\partial_1(y)$ con $y\in Z'_1$. Por la conmutatividad,
\[
  \phi(x)=\phi(\partial_1(y))=\partial_2(\phi'(y))\in\operatorname{Im}(\partial_2).
\]
Con un argumento completamente análogo, $\bar\psi$ está bien definida.

\paragraph{Exactitud en $\ker\partial_2$.} Claramente $\psi''\circ\phi''=0$, porque son restricciones de $\psi'$ y $\phi'$, y la composición de estas ya es cero (la primera fila es exacta). Por lo tanto, $\operatorname{Im}\phi''\subseteq\ker\psi''$.

Para la otra inclusión, sea $\phi''(a)\in\operatorname{Im}\phi''\subseteq\operatorname{Im}\phi'$. Como $\operatorname{Im}\phi'=\ker\psi'$, se tiene $\phi''(a)\in\ker\psi'$. Por otro lado, $\phi''(a)\in\ker\partial_2$, porque $\phi''$ está bien definida. Así, el elemento vive tanto en $\ker\psi'$ como en $\ker\partial_2$, es decir, en $\ker\psi''$.

\medskip\noindent\textbf{Ejercicio.} Demostrar la exactitud en $Z_2/\operatorname{Im}\partial_2$.

\medskip
Esta es la parte sencilla de la demostración. Ahora hay que definir el homomorfismo $\delta_*$, pues el enunciado solo dice que existe.

\paragraph{Definición de $\delta_*$.} La idea es viajar por el diagrama; esto es lo que se conoce como \emph{cacería de diagramas}. Sea $c'_3\in\ker(\partial_3)$; en particular, $c'_3\in Z'_3$, porque $\ker\partial_3$ es un submódulo de $Z'_3$. Como $\psi'$ es suprayectivo (por el $0$ al final de la primera fila), existe $c'_2\in Z'_2$ tal que $\psi'(c'_2)=c'_3$. Luego se baja este elemento con $\partial_2$ y se obtiene $\partial_2(c'_2)\in Z_2$. Se quiere ver que $\partial_2(c'_2)$ está en el núcleo de $\psi$, para poder jalarlo a $Z_1$. Por la conmutatividad,
\[
  \psi(\partial_2(c'_2))=\partial_3(\psi'(c'_2))=\partial_3(c'_3)=0,
\]
pues $c'_3\in\ker\partial_3$. Como $\ker\psi=\operatorname{Im}\phi$ y $\phi$ es inyectivo (por el $0$ al principio de la segunda fila), existe un único $c_1\in Z_1$ tal que
\[
  \phi(c_1)=\partial_2(c'_2).
\]
Recapitulando: se tomó el elemento $c'_3$, se jaló con la suprayectividad de $\psi'$, se empujó con $\partial_2$, se volvió a jalar con la inyectividad de $\phi$ (verificando que el elemento estaba en el núcleo de $\psi$), y ahora se empuja al cociente. Así que se define
\[
  \delta_*(c'_3)=c_1+\operatorname{Im}\partial_1 .
\]

\paragraph{$\delta_*$ está bien definido.} Hay un pequeño problema: al construir $\delta_*(c'_3)$ se tomó una preimagen $c'_2$ de $c'_3$, que existe por la suprayectividad de $\psi'$ pero no es única; hay una elección. Empujar no causa problemas, y jalar con $\phi$ tampoco, por la inyectividad; lo único que causa problemas es esa elección. Sea $c''_2\in(\psi')^{-1}(c'_3)$ otra preimagen. Entonces
\[
  c'_2-c''_2\in\ker\psi'=\operatorname{Im}(\phi')\ \Longrightarrow\ \partial_2(c'_2)-\partial_2(c''_2)\in\partial_2(\operatorname{Im}\phi')=\phi\,\partial_1(Z'_1),
\]
donde la igualdad del final es por la conmutatividad. Por lo tanto, la preimagen de $\partial_2(c'_2)-\partial_2(c''_2)$ en $Z_1$ vive en $\operatorname{Im}\partial_1$, lo cual implica que $\delta_*(c'_3)$ no depende de la elección. Este es el chiste del álgebra homológica: hay que ver el diagrama y perseguir los elementos hasta conseguir lo que se quiere.

\paragraph{Exactitud en $\ker\partial_3$.} Primero se ve que $\delta_*\circ\psi''=0$, es decir, que $\operatorname{Im}\psi''\subseteq\ker\delta_*$. Sea $x\in\ker\partial_2$. Para calcular $\delta_*(\psi''(x))$, se ve $\psi''(x)$ en $Z'_3$ y se toma una preimagen bajo $\psi'$; pero esa preimagen se puede tomar justamente como el elemento $x$. Luego $x$ se empuja con $\partial_2$ y se jala a $Z_1$:
\[
  \delta_*(\psi''(x))=\phi^{-1}(\partial_2(x))+\operatorname{Im}\partial_1=0,
\]
pues $x\in\ker\partial_2$, así que $\partial_2(x)=0$.

Ahora la otra inclusión, $\ker\delta_*\subseteq\operatorname{Im}\psi''$. Sea $c'_3\in\ker(\delta_*)$, y sean $c'_2$ y $c_1$ los elementos que se obtienen al bajar por el diagrama como en la definición de $\delta_*$. Que $\delta_*(c'_3)=0$ quiere decir que $c_1$ vive en la imagen de $\partial_1$: $c_1=\partial_1(c'_1)$ para algún $c'_1\in Z'_1$. Así que
\[
  \partial_2(c'_2)=\phi(c_1)=\phi(\partial_1(c'_1))=\partial_2(\phi'(c'_1)),
\]
donde la última igualdad es por la conmutatividad. Esto implica que $c'_2-\phi'(c'_1)\in\ker\partial_2$. Entonces, restando y sumando adentro y usando que $\psi'\circ\phi'=0$ por la exactitud de la primera fila,
\[
  c'_3=\psi'(c'_2)=\psi'(c'_2-\phi'(c'_1))+\psi'(\phi'(c'_1))=\psi'(c'_2-\phi'(c'_1))=\psi''(c'_2-\phi'(c'_1)),
\]
pues el elemento dentro de los paréntesis vive en $\ker\partial_2$, y entonces se le puede poner otra prima a la $\psi$. Así que $c'_3\in\operatorname{Im}\psi''$, y por lo tanto $\ker\delta_*=\operatorname{Im}\psi''$.

\paragraph{Exactitud en $Z_1/\operatorname{Im}\partial_1$.} Primero se ve que $\bar\phi\circ\delta_*=0$. Sea $c'_3\in\ker\partial_3$, con $c'_2$ y $c_1$ como en la definición de $\delta_*$. Entonces
\[
  \bar\phi(\delta_*(c'_3))=\phi(c_1)+\operatorname{Im}\partial_2=\partial_2(c'_2)+\operatorname{Im}\partial_2,
\]
el cual es cero en $Z_2/\operatorname{Im}\partial_2$. Así, $\operatorname{Im}\delta_*\subseteq\ker\bar\phi$.

Para la otra inclusión, sea $c_1+\operatorname{Im}\partial_1\in\ker(\bar\phi)$. Entonces $\phi(c_1)\in\operatorname{Im}\partial_2$, es decir, existe $c'_2\in Z'_2$ tal que $\phi(c_1)=\partial_2(c'_2)$. Sea $c'_3=\psi'(c'_2)$. Este elemento vive en $\ker\partial_3$:
\[
  \partial_3(c'_3)=\partial_3(\psi'(c'_2))=\psi(\partial_2(c'_2))=\psi(\phi(c_1))=0,
\]
por la conmutatividad y la exactitud de la segunda fila. Además, por construcción, $\delta_*(c'_3)=c_1+\operatorname{Im}\partial_1$: el elemento $c'_3$ se construyó justamente para que, al bajarlo y regresarlo, se cayera en el $c_1$ que se tomó desde el principio. Así que $c_1+\operatorname{Im}\partial_1\in\operatorname{Im}(\delta_*)$.

\medskip
Con esto queda demostrado que todo está bien definido y que la sucesión, conectada por la serpiente, es exacta. \qed

\medskip
Es una demostración larga, pero muy rica, porque hay que verificar muchas cosas. Ya con la demostración hecha, uno se puede quedar simplemente con el enunciado del lema de la serpiente, que se usará en el siguiente video para demostrar el teorema fundamental del álgebra homológica.

\begin{erratas}
  \errata{311}{5:11}{Se dice \guillemotleft{}análogamente $\bar\phi$ está bien definida\guillemotright{}, pero debe decir $\bar\psi$.}
  \errata{343}{5:43}{Se escribe $\psi''\phi''=0\Rightarrow\ker\psi''\subseteq\operatorname{Im}\phi''$, pero la inclusión que se obtiene es $\operatorname{Im}\phi''\subseteq\ker\psi''$.}
  \errata{352}{5:52}{El argumento de \guillemotleft{}la otra inclusión\guillemotright{} parte de un elemento $\phi''(a)\in\operatorname{Im}\phi''$ y concluye que está en $\ker\psi''$, es decir, vuelve a probar $\operatorname{Im}\phi''\subseteq\ker\psi''$. La inclusión que falta, $\ker\psi''\subseteq\operatorname{Im}\phi''$, se prueba así: si $b\in\ker\partial_2$ y $\psi'(b)=0$, por la exactitud de la primera fila $b=\phi'(a)$ con $a\in Z'_1$; como $\phi(\partial_1 a)=\partial_2(\phi'(a))=\partial_2 b=0$ y $\phi$ es inyectivo, $a\in\ker\partial_1$, y entonces $b=\phi''(a)$.}
  \errata{579}{9:39}{Se dice \guillemotleft{}como $\phi$ es inyectiva, existe un único $c_1$\guillemotright{}, pero la existencia de $c_1$ se sigue de que $\ker\psi=\operatorname{Im}\phi$ (exactitud de la segunda fila); la inyectividad de $\phi$ da la unicidad.}
  \errata{614}{10:14}{Se define $\delta_*(c'_3)=c'_2+\operatorname{Im}\partial_2$ (a las 11:43 se corrige el índice a $\partial_1$), pero debe decir $\delta_*(c'_3)=c_1+\operatorname{Im}\partial_1$.}
\end{erratas}

\end{document}
